\documentclass[journal]{IEEEtran}
\usepackage{amsmath,amssymb,amsfonts,amsthm}
\usepackage{algorithmic}
\usepackage{graphicx}
\usepackage{textcomp}
\usepackage{xcolor}
\definecolor{revisionred}{RGB}{210,38,48}
\newcommand{\rev}[1]{{\color{revisionred}#1}}
\usepackage{url}

% Theorem environments
\newtheorem{definition}{Definition}
\newtheorem{theorem}{Theorem}
\newtheorem{lemma}{Lemma}
\newtheorem{corollary}{Corollary}
\newtheorem{remark}{Remark}
\newtheorem{assumption}{Assumption}

% Standalone appendix numbering
\renewcommand{\theequation}{A.\arabic{equation}}
\renewcommand{\thefigure}{A.\arabic{figure}}
\renewcommand{\thedefinition}{A.\arabic{definition}}
\renewcommand{\thetheorem}{A.\arabic{theorem}}
\renewcommand{\thelemma}{A.\arabic{lemma}}
\renewcommand{\thecorollary}{A.\arabic{corollary}}
\renewcommand{\theremark}{A.\arabic{remark}}
\renewcommand{\theassumption}{A.\arabic{assumption}}

\onecolumn
\allowdisplaybreaks

\begin{document}
\color{revisionred}

\title{Technical Appendix: Convergence of Fixed-Participation, Dynamic-Range-Aware Top-$k$ OFDM-AirComp Federated Learning}
\author{Jinning Song and Dong Liu}
\maketitle

\section{Algorithm Flow and Assumptions}

To facilitate the convergence analysis, we first formalize the algorithm flow, the fixed-participation OFDM-AirComp aggregation and receiver dynamic-range model, and the assumptions used in the proof. Throughout this appendix, the global objective is
\begin{equation}
    f(\theta) \triangleq \frac{1}{N}\sum_{i=1}^{N} f_i(\theta),
\end{equation}
where $\theta\in\mathbb{R}^d$ is the global model parameter. We use a mechanism index
\begin{equation}
    A\in\{\mathrm{Top}\text{-}k,\mathrm{Rand}\text{-}k,\mathrm{Full}\}
\end{equation}
to unify the proposed top-$k$ mechanism and the baseline mechanisms. The proposed method corresponds to $A=\mathrm{Top}\text{-}k$. The full-update case is recovered by setting $k=d$ and $\mathcal{S}_{\mathrm{Full}}^d(v)=v$.

The convergence bound is stated for a generic admissible mechanism $A$. The mechanism-specific learning advantage is captured by a certified energy-retention contraction coefficient $\bar{\omega}_A(k)$ that is valid along the training trajectory, while the mechanism-specific hardware effect is captured by the dynamic-range residual functional $E_{\mathrm{clip},A}^t(k,b_t)$ and its support-overlap interpretation. In the proof, $\bar{\omega}_A(k)$ is used as a uniform lower bound, not merely as an a posteriori time-average statistic. Empirical retained-energy averages can be used in the main-text design as calibration surrogates, but the theorem itself relies on the certified contraction condition stated in Assumption~\ref{assum:contraction}.

\subsection{Algorithm Flow and Error Feedback Mechanism}

At the beginning of round $t$, client $i$ initializes its local model as
\begin{equation}
    \theta_i^{t,0}=\theta^t.
\end{equation}
It then performs $\tau$ local SGD steps
\begin{equation}
    \theta_i^{t,s}=\theta_i^{t,s-1}-\eta g_i^{t,s-1},
    \qquad s=1,\ldots,\tau,
\end{equation}
where $g_i^{t,s-1}$ denotes the stochastic gradient evaluated on the local mini-batch at step $s$. The raw accumulated local update is
\begin{equation}
    \Delta_i^t\triangleq u_i^t
    =\theta_i^{t,\tau}-\theta^t
    =-\eta\sum_{s=1}^{\tau}g_i^{t,s-1}.
\end{equation}

Each client maintains a local error memory $e_i^t$, initialized as $e_i^{-1}=\mathbf{0}$. The sparsification-then-clipping mechanism proceeds as follows:
\begin{enumerate}
    \item \textbf{Compensation:} form the compensated update
    \begin{equation}
        v_i^t=\Delta_i^t+e_i^{t-1}.
    \end{equation}
    \item \textbf{Mechanism-dependent sparsification:} apply
    \begin{equation}
        \tilde v_{A,i}^t(k)=\mathcal{S}_A^k(v_i^t).
    \end{equation}
    For $A=\mathrm{Top}\text{-}k$, $\mathcal{S}_A^k$ keeps the $k$ largest-magnitude entries. For $A=\mathrm{Rand}\text{-}k$, one uniformly sampled support is shared by all clients within a round and resampled independently across rounds. For $A=\mathrm{Full}$, $k=d$ and no sparsification is applied.
    \item \textbf{Error-memory update:} store only the sparsification residual
    \begin{equation}
        e_i^t=v_i^t-\tilde v_{A,i}^t(k).
    \end{equation}
    \item \textbf{Pre-transmission clipping:} clip the sparse update element-wise as
    \begin{equation}
        s_{A,i}^t(k)=\mathrm{clip}\big(\tilde v_{A,i}^t(k),\eta\tau C\big).
    \end{equation}
\end{enumerate}
The clipping step guarantees
\begin{equation}
    \|s_{A,i}^t(k)\|_0\le k,
    \qquad
    \|s_{A,i}^t(k)\|_{\infty}\le \eta\tau C.
\end{equation}
We define the pre-transmission clipping bias as
\begin{equation}
    \xi_{A,i}^t(k)\triangleq \tilde v_{A,i}^t(k)-s_{A,i}^t(k).
\end{equation}
Then the transmitted signal admits the exact decomposition
\begin{equation}
    \label{eqn:transmit_signal}
    s_{A,i}^t(k)
    =\Delta_i^t+e_i^{t-1}-e_i^t-\xi_{A,i}^t(k).
\end{equation}
The clipping residual is not fed back into $e_i^t$, which prevents clipping distortion from being recursively accumulated by the error-feedback memory.

\subsection{Fixed-Participation OFDM-AirComp Aggregation and Receiver Dynamic-Range Model}

Every logical communication round is completed by all $N$ clients; there is no algorithmic client sampling and the learning participation set is fixed. A candidate burst is counted only if $|h_i^t|\ge h_{\mathrm{cut}}$ for every $i=1,\ldots,N$; otherwise the physical layer retries the same client set. Such waiting affects communication latency but does not change the aggregate, the denominator, the error memories, or the round index used in the convergence analysis. Equivalently, the analysis is conditioned on successful full-client alignment in every counted round.

The real-valued sparse vector $s_{A,i}^t(k)$ is mapped onto $S=\lceil d/M\rceil$ OFDM symbols, each containing $M$ data subcarriers. Let $S_{A,i,\ell}^t[m;k]$ denote the real model coordinate assigned to the $m$-th subcarrier of the $\ell$-th OFDM symbol. The client applies complex channel inversion
\begin{equation}
    X_{A,i,\ell}^t[m]
    =\frac{b_t}{g_i^t}S_{A,i,\ell}^t[m;k],
    \qquad i=1,\ldots,N,
\end{equation}
where $g_i^t=\sqrt{\beta_i}h_i^t$ is the public frequency-flat channel coefficient in round $t$. Conditioned on successful alignment, the frequency-domain observation is
\begin{equation}
    Y_{A,\ell}^t[m;k,b_t]
    =b_t\sum_{i=1}^{N}S_{A,i,\ell}^t[m;k]+N_\ell^t[m],
\end{equation}
with $N_\ell^t[m]\sim\mathcal{CN}(0,\sigma_{\mathrm{sc}}^2)$. The real part of this noise is the Gaussian randomness used for the per-round client-level privacy guarantee.

Let $c_{\mathrm{tx}}\triangleq\eta\tau C$. The complete-burst average transmit power of client $i$ is
\begin{equation}
    P_{A,i}^{t}
    =\frac{1}{SM}\sum_{\ell=1}^{S}\sum_{m=0}^{M-1}|X_{A,i,\ell}^t[m]|^2
    =\frac{b_t^2}{SM|g_i^t|^2}\|s_{A,i}^t(k)\|_2^2.
\end{equation}
Since $\|s_{A,i}^t(k)\|_0\le k$ and $\|s_{A,i}^t(k)\|_\infty\le c_{\mathrm{tx}}$, the public per-round power ceiling is
\begin{equation}
    \label{eqn:power_ceiling}
    B_P^t(k)
    \triangleq
    \min_{1\le i\le N}
    \frac{|g_i^t|\sqrt{SM P_{\mathrm{cap}}}}
    {c_{\mathrm{tx}}\sqrt{k}}.
\end{equation}
For replace-one client adjacency, the sparse-query sensitivity before multiplication by $b_t$ is
\begin{equation}
    \Delta(k)=2c_{\mathrm{tx}}\sqrt{k}.
\end{equation}
Define the exact sufficient per-round privacy ceiling
\begin{equation}
    \label{eqn:per_round_privacy_ceiling}
    B_{\epsilon,\mathrm{ex}}(k)
    \triangleq
    \frac{\sigma_{\mathrm{sc}}}{\Delta(k)}
    \left(
    \sqrt{\epsilon+\ln(1/\delta)}-\sqrt{\ln(1/\delta)}
    \right).
\end{equation}
The protocol uses the unique per-round scaling rule
\begin{equation}
    \label{eqn:per_round_scaling}
    b_t=b_t^\star(k)
    \triangleq
    \min\left\{B_{\epsilon,\mathrm{ex}}(k),B_P^t(k)\right\}.
\end{equation}
Thus $b_t$ may vary across rounds through the current public channel state, but it does not depend on current private update norms and requires no cross-round privacy-budget allocation.

Let $Q=L_{\mathrm{os}}M$ and let $\mathcal Z_{\mathrm{os}}$ denote centered zero padding. The power-preserving oversampled vector is
\begin{equation}
    \mathbf Y_{A,\ell,\mathrm{os}}^t[k,b_t]
    =\sqrt{Q/M}\,\mathcal Z_{\mathrm{os}}\big(\mathbf Y_{A,\ell}^t[k,b_t]\big),
\end{equation}
and the unitary-IFFT receiver waveform is
\begin{equation}
    y_{A,\ell}^t[n;k,b_t]
    =\mathrm{IFFT}_{Q}^{\mathrm u}\!\left(\mathbf Y_{A,\ell,\mathrm{os}}^t[k,b_t]\right)[n].
\end{equation}
The ideal per-round RMS-AGC defines
\begin{equation}
    P_{\mathrm{avg}}^t
    =\frac{1}{SQ}\sum_{\ell=1}^{S}\sum_{n=0}^{Q-1}|y_{A,\ell}^t[n;k,b_t]|^2,
    \quad A_{\mathrm{rms}}^t=\sqrt{P_{\mathrm{avg}}^t},
    \quad a_t=(A_{\mathrm{rms}}^t)^{-1}.
\end{equation}
The gain is fixed within the round and is explicitly reversed before model-update normalization. The finite complex-envelope dynamic range is modeled by the phase-preserving radial limiter
\begin{equation}
    \mathcal C_\gamma(z)=
    \begin{cases}
      z,& |z|\le\gamma,\\
      \gamma z/|z|,& |z|>\gamma,
    \end{cases}
    \qquad \mathcal C_\gamma(0)=0.
\end{equation}
The input-referred residual is
\begin{equation}
    r_{\mathrm{clip},A,\ell}^t[n;k,b_t]
    =a_t^{-1}\mathcal C_\gamma\!\left(a_ty_{A,\ell}^t[n;k,b_t]\right)
    -y_{A,\ell}^t[n;k,b_t],
\end{equation}
with absolute threshold $A_{\max}^t=\gamma A_{\mathrm{rms}}^t$ and exact residual energy
\begin{equation}
    \label{eqn:eclip_def}
    E_{\mathrm{clip},A}^t(k,b_t)
    =\frac{M}{Q}\sum_{\ell=1}^{S}\sum_{n=0}^{Q-1}
    \left(|y_{A,\ell}^t[n;k,b_t]|-A_{\max}^t\right)_+^2.
\end{equation}
After the unitary FFT, extraction of the original bins, and inverse oversampling normalization, let $\mathbf R_{\mathrm{clip},A,\mathcal D}^t$ be the concatenated frequency-domain residual. The equivalent receiver dynamic-range aggregation error and channel-noise error are
\begin{equation}
    e_{\mathrm{dr},A}^t
    =\frac{\mathrm{Re}\{\mathbf R_{\mathrm{clip},A,\mathcal D}^t\}}{b_tN},
    \qquad
    e_{\mathrm{ch}}^t
    =\frac{\mathrm{Re}\{\mathbf N_{\mathcal D}^t\}}{b_tN},
\end{equation}
where
\begin{equation}
    \mathbb E\|e_{\mathrm{ch}}^t\|_2^2
    =\frac{d\sigma_{\mathrm{sc}}^2}{2b_t^2N^2}.
\end{equation}
The physical update is therefore
\begin{equation}
    \label{eqn:physical_update_dr}
    \theta^{t+1}
    =\theta^t+\frac1N\sum_{i=1}^{N}s_{A,i}^t(k)
    +e_{\mathrm{ch}}^t+e_{\mathrm{dr},A}^t.
\end{equation}

Let $\mathcal F_t$ denote the filtration generated by the physical iterate, all error memories, and all randomness before round $t$. The current public channel state and scaling are measurable before transmission. Channel noise is zero mean conditional on this enlarged public history; the radial-limiting residual is not assumed to be zero mean or independent.

\subsection{Core Assumptions}

\begin{assumption}[\bf Smoothness and Lower Bound]
    \label{assum:smoothness}
    Each local objective $f_i$ is $L$-smooth, and the global objective satisfies $f(\theta)\ge f^*>-\infty$ for all $\theta\in\mathbb{R}^d$.
\end{assumption}

\begin{assumption}[\bf Unbiased Stochastic Gradient]
    \label{assum:unbiased}
    Let $\mathcal{H}_{i}^{t,s-1}$ denote the local history before sampling the mini-batch at local step $s$ of round $t$. The stochastic gradient satisfies
    \begin{equation}
        \mathbb{E}\left[g_i^{t,s-1}\mid\mathcal{H}_{i}^{t,s-1}\right]
        =\nabla f_i(\theta_i^{t,s-1}).
    \end{equation}
\end{assumption}

\begin{assumption}[\bf Bounded Gradient Variance]
    \label{assum:variance}
    The stochastic gradient variance satisfies
    \begin{equation}
        \mathbb{E}\left[
        \|g_i^{t,s-1}-\nabla f_i(\theta_i^{t,s-1})\|_2^2
        \mid\mathcal{H}_{i}^{t,s-1}
        \right]\le \zeta^2.
    \end{equation}
\end{assumption}

\begin{assumption}[\bf Bounded Heterogeneity]
    \label{assum:heterogeneity}
    There exist constants $G^2\ge 1$ and $\kappa^2\ge0$ such that
    \begin{equation}
        \frac1N\sum_{i=1}^{N}\|\nabla f_i(\theta)\|_2^2
        \le G^2\|\nabla f(\theta)\|_2^2+\kappa^2.
    \end{equation}
\end{assumption}

\begin{assumption}[\bf Bounded Pre-Transmission Clipping Bias]
    \label{assum:clipping}
    The clipping bias satisfies
    \begin{equation}
        \mathbb{E}_t\|\xi_{A,i}^t(k)\|_2^2
        \le \eta^2\tau^2\beta_A^2(k),
        \qquad \forall i,t,
    \end{equation}
    where $\beta_A(k)$ depends on the clipping threshold, the sparsification mechanism, and the data distribution, but is independent of $t$.
\end{assumption}

\begin{assumption}[\bf Certified Energy-Retention Contraction]
    \label{assum:contraction}
    For each admissible mechanism $A$, there exists a certified contraction coefficient $\bar{\omega}_A(k)\in(0,1]$ such that, for all clients and all rounds along the training trajectory,
    \begin{equation}
        \mathbb{E}_t\|v_i^t-\mathcal{S}_A^k(v_i^t)\|_2^2
        \le
        \big(1-\bar{\omega}_A(k)\big)\mathbb{E}_t\|v_i^t\|_2^2,
        \qquad \forall i,t.
    \end{equation}
    Thus $\bar{\omega}_A(k)$ is a trajectory-wise lower bound used to close the error-feedback recursion. It should not be understood as only the empirical time-average retained-energy ratio. In implementation, one may estimate an empirical retained-energy statistic for design, but the proof requires the above uniform contraction condition or a certified lower envelope of the per-round coefficients.

    For uniform Rand-$k$ without rescaling, $\bar{\omega}_{\mathrm{Rand}}(k)=k/d$ under uniform coordinate sampling. For Top-$k$,
    \begin{equation}
        \frac{\|\mathrm{Top}_k(v)\|_2^2}{\|v\|_2^2}\ge \frac{k}{d},
    \end{equation}
    and $\bar{\omega}_{\mathrm{Top}}(k)$ can be larger than $k/d$ when the update energy is concentrated on important coordinates. For the full-update mechanism, $\bar{\omega}_{\mathrm{Full}}(d)=1$.
\end{assumption}

\begin{assumption}[\bf Receiver Dynamic-Range Residual]
    \label{assum:dr_residual}
    The expected radial-limiting residual energy is finite:
    \begin{equation}
        \mathbb E[E_{\mathrm{clip},A}^t(k,b_t)]<\infty.
    \end{equation}
    The aggregation error $e_{\mathrm{dr},A}^t$ is not assumed to be zero mean and is not assumed to be independent of the channel noise or local updates.
\end{assumption}

\begin{assumption}[\bf Successful Full-Client Alignment and Public Scaling]
    \label{assum:public_scaling}
    Every counted logical round satisfies $|h_i^t|\ge h_{\mathrm{cut}}$ for all clients and is completed by all $N$ clients under exact pre-equalized alignment. The public scaling is given by~\eqref{eqn:per_round_scaling}, is independent of the realized current private updates, and satisfies
    \begin{equation}
        \mathbb E\left[\frac{1}{b_t^2}\right]<\infty.
    \end{equation}
\end{assumption}

\begin{theorem}[\bf Per-Round Client-Level Privacy]
    \label{thm:per_round_privacy}
    Consider replace-one client adjacency. Conditional on the same preceding public transcript, current public channel state, public $b_t$, and data-independent randomness of unchanged clients, the noisy AirComp observation in round $t$ satisfies client-level $(\epsilon,\delta)$-DP whenever $0<b_t\le B_{\epsilon,\mathrm{ex}}(k)$. In particular, the protocol scaling~\eqref{eqn:per_round_scaling} satisfies this condition in every round.
\end{theorem}
\begin{proof}
    The pre-noise query has sensitivity at most $b_t\Delta(k)$. The real part of $\mathcal{CN}(0,\sigma_{\mathrm{sc}}^2)$ noise has variance $\sigma_{\mathrm{sc}}^2/2$. The order-$\alpha$ Gaussian RDP cost is therefore at most $\alpha b_t^2\Delta^2(k)/\sigma_{\mathrm{sc}}^2$. Applying the standard RDP-to-DP conversion and optimizing over $\alpha>1$ yields~\eqref{eqn:per_round_privacy_ceiling}. AGC, radial limiting, inverse-gain restoration, FFT recovery, division by $N$, and model updating are post-processing of the noisy observation.
\end{proof}

The theorem is a per-round statement. The paper does not claim that the complete $T$-round transcript satisfies the same $(\epsilon,\delta)$ without an additional composition accountant.

\section{Auxiliary Lemmas}

\rev{This standalone technical appendix uses independent A.x numbering. Its effective-contraction, cumulative-memory, local-divergence, and dynamic-range-error lemmas provide the detailed derivation behind the corresponding proof construction in the main text; Theorem~\ref{thm:main_convergence} is the appendix counterpart of the main dynamic-range-aware convergence theorem.}

\begin{lemma}[\bf Effective Contraction Step]
    \label{lemma:topk}
    Under Assumption~\ref{assum:contraction}, for every client $i$ and round $t$,
    \begin{equation}
        \mathbb{E}_t\|e_i^t\|_2^2
        \le
        \big(1-\bar{\omega}_A(k)\big)
        \mathbb{E}_t\|\Delta_i^t+e_i^{t-1}\|_2^2.
    \end{equation}
\end{lemma}
\begin{proof}
    Since $e_i^t=v_i^t-\mathcal{S}_A^k(v_i^t)$ and $v_i^t=\Delta_i^t+e_i^{t-1}$, the claim follows directly from Assumption~\ref{assum:contraction}.
\end{proof}

\begin{lemma}[\bf Cumulative Error Memory Bound]
    \label{lemma:error_bound}
    Under Assumption~\ref{assum:contraction}, the cumulative error memory satisfies
    \begin{equation}
        \sum_{t=0}^{T-1}\frac1N\sum_{i=1}^{N}\mathbb{E}\|e_i^t\|_2^2
        \le
        \mu_A(k)
        \sum_{t=0}^{T-1}\frac1N\sum_{i=1}^{N}\mathbb{E}\|\Delta_i^t\|_2^2,
    \end{equation}
    where, for $0<\bar{\omega}_A(k)<1$,
    \begin{equation}
        \label{eqn:mu_def}
        \mu_A(k)
        =\frac{2(1-\bar{\omega}_A(k))(2-\bar{\omega}_A(k))}{\bar{\omega}_A^2(k)}.
    \end{equation}
    For $\bar{\omega}_A(k)=1$, the compression residual vanishes and $\mu_A(k)=0$.
\end{lemma}
\begin{proof}
    The full-update case is immediate. For $0<\bar{\omega}_A(k)<1$, Lemma~\ref{lemma:topk} and Young's inequality imply that, for any $\alpha>0$,
    \begin{align}
        \mathbb{E}\|e_i^t\|_2^2
        &\le
        \big(1-\bar{\omega}_A(k)\big)(1+\alpha)\mathbb{E}\|e_i^{t-1}\|_2^2
        \notag\\
        &\quad+
        \big(1-\bar{\omega}_A(k)\big)(1+\alpha^{-1})\mathbb{E}\|\Delta_i^t\|_2^2.
    \end{align}
    Choose
    \begin{equation}
        \alpha=\frac{\bar{\omega}_A(k)}{2(1-\bar{\omega}_A(k))}.
    \end{equation}
    Then the recursion coefficient is
    \begin{equation}
        \varrho_A(k)
        =\big(1-\bar{\omega}_A(k)\big)(1+\alpha)
        =1-\frac{\bar{\omega}_A(k)}{2}<1,
    \end{equation}
    and the innovation coefficient is
    \begin{equation}
        \beta_{\mathrm{err},A}(k)
        =\big(1-\bar{\omega}_A(k)\big)(1+\alpha^{-1})
        =\frac{(1-\bar{\omega}_A(k))(2-\bar{\omega}_A(k))}{\bar{\omega}_A(k)}.
    \end{equation}
    Since $e_i^{-1}=\mathbf{0}$, unrolling the recursion and summing over $t$ give
    \begin{equation}
        \sum_{t=0}^{T-1}\mathbb{E}\|e_i^t\|_2^2
        \le
        \frac{\beta_{\mathrm{err},A}(k)}{1-\varrho_A(k)}
        \sum_{t=0}^{T-1}\mathbb{E}\|\Delta_i^t\|_2^2
        =\mu_A(k)\sum_{t=0}^{T-1}\mathbb{E}\|\Delta_i^t\|_2^2.
    \end{equation}
    Averaging over clients completes the proof.
\end{proof}

\begin{lemma}[\bf Bounded Local Divergence]
    \label{lemma:local_div}
    Under Assumptions~\ref{assum:smoothness}--\ref{assum:heterogeneity}, if $\tau\ge2$ and $\eta\le 1/(2\sqrt{2}\tau L)$, then for every local step $s\le\tau$,
    \begin{align}
        \frac1N\sum_{i=1}^{N}\mathbb{E}_t\|\theta^t-\theta_i^{t,s-1}\|_2^2
        &\le
        16\eta^2\tau^2G^2\mathbb{E}_t\|\nabla f(\theta^t)\|_2^2
        \notag\\
        &\quad+16\eta^2\tau^2\kappa^2+4\tau\eta^2\zeta^2.
    \end{align}
\end{lemma}
\begin{proof}
    Define $d_i^{t,p}=\theta_i^{t,p}-\theta^t$. Then $d_i^{t,0}=\mathbf{0}$ and
    \begin{equation}
        d_i^{t,p}=d_i^{t,p-1}-\eta g_i^{t,p-1}.
    \end{equation}
    Decompose $g_i^{t,p-1}=\nabla f_i(\theta_i^{t,p-1})+\varepsilon_i^{t,p-1}$, where the noise term is conditionally zero-mean and has conditional variance bounded by $\zeta^2$. Following the standard local-SGD drift argument, Young's inequality with parameter $(2\tau-1)^{-1}$ gives
    \begin{align}
        \mathbb{E}_t\|d_i^{t,p}\|_2^2
        &\le
        \left(1+\frac{1}{2\tau-1}+4\eta^2L^2\tau\right)
        \mathbb{E}_t\|d_i^{t,p-1}\|_2^2
        \notag\\
        &\quad+4\eta^2\tau\mathbb{E}_t\|\nabla f_i(\theta^t)\|_2^2+\eta^2\zeta^2.
    \end{align}
    Averaging over clients, using Assumption~\ref{assum:heterogeneity}, and applying $\eta\le 1/(2\sqrt{2}\tau L)$ yield
    \begin{align}
        \frac1N\sum_{i=1}^{N}\mathbb{E}_t\|d_i^{t,p}\|_2^2
        &\le
        \left(1+\frac{1}{\tau-1}\right)
        \frac1N\sum_{i=1}^{N}\mathbb{E}_t\|d_i^{t,p-1}\|_2^2
        \notag\\
        &\quad+4\eta^2\tau G^2\mathbb{E}_t\|\nabla f(\theta^t)\|_2^2
        +4\eta^2\tau\kappa^2+\eta^2\zeta^2.
    \end{align}
    Unrolling for at most $\tau-1$ steps and using $(1+1/(\tau-1))^{\tau}\le4$ gives the result.
\end{proof}

\begin{lemma}[\bf Receiver Dynamic-Range Aggregation Error Bound]
    \label{lemma:dr_error}
    Under the unitary FFT/IFFT operations and the radial-limiting residual definition, the equivalent receiver dynamic-range aggregation error satisfies the pathwise inequality
    \begin{equation}
        \|e_{\mathrm{dr},A}^t\|_2^2
        \le
        \frac{E_{\mathrm{clip},A}^t(k,b_t)}{b_t^2N^2}.
    \end{equation}
    Consequently,
    \begin{equation}
        \mathbb{E}_t\|e_{\mathrm{dr},A}^t\|_2^2
        \le
        \frac{\mathbb{E}_t[E_{\mathrm{clip},A}^t(k,b_t)]}{b_t^2N^2}.
    \end{equation}
\end{lemma}
\begin{proof}
    By definition,
    \begin{equation}
        \|e_{\mathrm{dr},A}^t\|_2^2
        =\frac{\|\mathrm{Re}\{\mathbf{R}_{\mathrm{clip},A,\mathcal{D}}^t\}\|_2^2}{b_t^2N^2}.
    \end{equation}
    Taking the real part and projecting onto data subcarriers cannot increase energy. Hence,
    \begin{equation}
        \|\mathrm{Re}\{\mathbf{R}_{\mathrm{clip},A,\mathcal{D}}^t\}\|_2^2
        \le
        \sum_{\ell=1}^{S}\|\mathbf{R}_{\mathrm{clip},A,\ell}^t\|_2^2.
    \end{equation}
    Parseval's theorem under unitary FFT/IFFT normalization and the inverse oversampling factor gives
    \begin{equation}
        \sum_{\ell=1}^{S}\|\mathbf{R}_{\mathrm{clip},A,\ell}^t\|_2^2
        \le\frac{M}{Q}\sum_{\ell=1}^{S}\|\mathbf{r}_{\mathrm{clip},A,\ell}^t\|_2^2
        =E_{\mathrm{clip},A}^t(k,b_t).
    \end{equation}
    Combining these inequalities proves the pathwise result. Taking conditional expectation gives the second claim.
\end{proof}

\begin{lemma}[\bf Selection-to-Peak Stress Bound]
    \label{lemma:peak_stress}
    Define the subcarrier concurrency
    \begin{equation}
        U_{A,\ell,m}^t(k)=\sum_{i=1}^{N}\mathbf{1}\{S_{A,i,\ell}^t[m;k]\neq0\},
    \end{equation}
    and the peak-stress quantity
    \begin{equation}
        \Omega_{A,\ell}^t(k)=\sum_{m=0}^{M-1}\big(U_{A,\ell,m}^t(k)\big)^2.
    \end{equation}
    Let
    \begin{equation}
        G_{A,\ell}^t[m;k]=\sum_{i=1}^{N}S_{A,i,\ell}^t[m;k]
    \end{equation}
    be the noiseless aggregated frequency-domain signal before multiplication by $b_t$. Then the noiseless time-domain signal satisfies
    \begin{equation}
        \max_n\left|y_{A,\ell,\mathrm{sig}}^t[n;k,b_t]\right|^2
        \le
        \frac{Q}{M}b_t^2\eta^2\tau^2C^2\Omega_{A,\ell}^t(k).
    \end{equation}
\end{lemma}
\begin{proof}
    Since $|S_{A,i,\ell}^t[m;k]|\le\eta\tau C$,
    \begin{equation}
        |G_{A,\ell}^t[m;k]|\le \eta\tau C\,U_{A,\ell,m}^t(k).
    \end{equation}
    Hence
    \begin{equation}
        \|\mathbf{G}_{A,\ell,\mathrm{os}}^t(k)\|_2^2
        =\frac{Q}{M}\|\mathbf{G}_{A,\ell}^t(k)\|_2^2
        \le \frac{Q}{M}\eta^2\tau^2C^2\Omega_{A,\ell}^t(k).
    \end{equation}
    A unitary IFFT preserves the Euclidean norm, and the maximum entry magnitude is upper-bounded by the vector norm. Multiplying by $b_t^2$ proves the result.
\end{proof}

\begin{lemma}[\bf Pairwise Support Overlap Identity]
    \label{lemma:overlap}
    Let $\mathcal{K}_{A,i}^t(k)$ be the selected support of client $i$ in round $t$, and let $\mathcal{I}_{A,i}^t(k)\subseteq\mathcal{K}_{A,i}^t(k)$ be the actually active nonzero support after sparsification and pre-transmission clipping. Define
    \begin{equation}
        \Omega_A^t(k)=\sum_{\ell=1}^{S}\Omega_{A,\ell}^t(k).
    \end{equation}
    Then
    \begin{align}
        \mathbb{E}[\Omega_A^t(k)]
        &=\sum_{i=1}^{N}\mathbb{E}\big[|\mathcal{I}_{A,i}^t(k)|\big]
        \notag\\
        &\quad+
        \sum_{i\neq j}\mathbb{E}\big[|\mathcal{I}_{A,i}^t(k)\cap\mathcal{I}_{A,j}^t(k)|\big].
    \end{align}
    If each selected support has cardinality $k$ and the active support coincides with the selected support, define
    \begin{equation}
        \rho_A^t(k)
        =\frac{1}{N(N-1)k}
        \sum_{i\neq j}\mathbb{E}\big[|\mathcal{K}_{A,i}^t(k)\cap\mathcal{K}_{A,j}^t(k)|\big].
    \end{equation}
    In this ideal exactly-$k$ case,
    \begin{equation}
        \mathbb{E}[\Omega_A^t(k)]
        =Nk+N(N-1)k\rho_A^t(k).
    \end{equation}
    More generally, when clipping or zero-valued selected entries reduce the active support so that $|\mathcal{I}_{A,i}^t(k)|\le k$, the same expression provides the upper bound
    \begin{equation}
        \mathbb{E}[\Omega_A^t(k)]
        \le Nk+N(N-1)k\rho_A^t(k),
    \end{equation}
    provided $\rho_A^t(k)$ is defined on the selected supports.
\end{lemma}
\begin{proof}
    Ignoring padded inactive entries, write
    \begin{equation}
        \Omega_A^t(k)=\sum_{m=1}^{d}\left(\sum_{i=1}^{N}q_{A,i,m}^t(k)\right)^2,
    \end{equation}
    where $q_{A,i,m}^t(k)=\mathbf{1}\{m\in\mathcal{I}_{A,i}^t(k)\}$. Expanding the square gives the self-support term and the pairwise overlap term. Taking expectations proves the first identity. If $\mathcal{I}_{A,i}^t(k)=\mathcal{K}_{A,i}^t(k)$ and $|\mathcal{K}_{A,i}^t(k)|=k$, substituting the definition of $\rho_A^t(k)$ gives the equality. If $\mathcal{I}_{A,i}^t(k)\subseteq\mathcal{K}_{A,i}^t(k)$, the active self-support and pairwise active overlaps are no larger than their selected-support counterparts, which gives the upper bound.
\end{proof}

\begin{corollary}[\bf High-Probability Selection-to-Clipping Linkage]
    \label{cor:clipping_linkage}
    Let the time-domain noise-envelope event be
    \begin{equation}
        \mathcal{N}_{\delta}^{t}
        =\left\{\max_{\ell,n}|y_{\ell,\mathrm{noise}}^t[n]|\le\nu_\delta\right\},
        \qquad
        \mathbb{P}(\mathcal{N}_{\delta}^{t})\ge 1-\delta.
    \end{equation}
    On $\mathcal{N}_{\delta}^{t}$,
    \begin{equation}
        E_{\mathrm{clip},A}^{t}(k,b_t)
        \le
        \sum_{\ell=1}^{S}M
        \left(
        \sqrt{\frac{Q}{M}}b_t\eta\tau C\sqrt{\Omega_{A,\ell}^{t}(k)}+
        \nu_\delta-A_{\max}^t
        \right)_+^2.
    \end{equation}
\end{corollary}
\begin{proof}
    On $\mathcal{N}_{\delta}^{t}$, Lemma~\ref{lemma:peak_stress} implies
    \begin{equation}
        |y_{A,\ell}^t[n;k,b_t]|
        \le
        \sqrt{\frac{Q}{M}}b_t\eta\tau C\sqrt{\Omega_{A,\ell}^{t}(k)}+\nu_\delta.
    \end{equation}
    The radial residual magnitude is $(|y|-A_{\max}^t)_+$. Summing over all $Q$ samples and applying the normalization $M/Q$ completes the proof.
\end{proof}

\begin{remark}[\bf Role of the Selection-Dependent Clipping Linkage]
    Corollary~\ref{cor:clipping_linkage} is not used to replace the exact dynamic-range residual in the main convergence theorem. The theorem keeps $\mathbb{E}[E_{\mathrm{clip},A}^{t}(k,b_t)]$ explicitly. The corollary explains how the choice of $A$ influences dynamic-range distortion through the support-overlap path
    \begin{equation}
        A\rightarrow \rho_A^t(k)\rightarrow\Omega_A^t(k)\rightarrow E_{\mathrm{clip},A}^{t}(k,b_t).
    \end{equation}
\end{remark}

\section{Proof of Main Convergence Bound}

\begin{theorem}[\bf Dynamic-Range-Aware Convergence Bound]
    \label{thm:main_convergence}
    Suppose Assumptions~\ref{assum:smoothness}--\ref{assum:public_scaling} hold, $\tau\ge2$, and the learning rate satisfies
    \begin{equation}
        \label{eqn:lr_condition_new}
        \eta
        \le
        \min\left\{
            \frac{1}{128L\tau G^2},
            \frac{1}{\tau GL\sqrt{1+384\mu_A(k)}}
        \right\},
    \end{equation}
    where $\mu_A(k)$ is defined in~\eqref{eqn:mu_def}. Define
    \begin{equation}
        \Gamma_A(k)=2+16L^2\mu_A(k)\eta^2\tau^2G^2,
    \end{equation}
    \begin{equation}
        C_{\mathrm{dr}}=\frac{4}{\eta\tau}+L,
    \end{equation}
    \begin{equation}
        \mathcal{E}_A(k)
        \triangleq
        (\eta\tau+2L\eta^2\tau^2)\beta_A^2(k),
    \end{equation}
    \begin{equation}
        \Lambda
        \triangleq
        \eta\tau\left(\frac12+4L\eta\tau\right)
        \left(1+32\eta^2\tau^2G^2L^2\right),
    \end{equation}
    \begin{equation}
        R_A(k)=8L^2\mu_A(k)\eta^2\tau^2(2\kappa^2+\zeta^2),
    \end{equation}
    and
    \begin{equation}
        \Psi
        \triangleq
        2(2\kappa^2+\zeta^2)
        +\eta L\left(\frac12+4L\eta\tau\right)
        (32\tau\kappa^2+8\zeta^2),
    \end{equation}
    which is zero when $\kappa^2+\zeta^2=0$. Let $\Delta f\triangleq f(\theta^0)-f^*$. Define the learning-side term
    \begin{align}
        \label{eqn:phi_def_theorem}
        \Phi_A\big(k;\bar{\omega}_A,\beta_A\big)
        &\triangleq
        \frac{16\Gamma_A(k)\Delta f}{T\eta\tau}
        +\frac{16\Gamma_A(k)\mathcal{E}_A(k)}{\eta\tau}
        +16\Gamma_A(k)\eta\tau L\Psi
        \notag\\
        &\quad+
        \frac{16\Gamma_A(k)\Lambda R_A(k)}{\eta\tau}
        +R_A(k).
    \end{align}
    Then the physical iterates generated by~\eqref{eqn:physical_update_dr} satisfy
    \begin{align}
        \label{eqn:final_dr_convergence_bound_theorem}
        \frac1T\sum_{t=0}^{T-1}\mathbb{E}\\|\nabla f(\theta^t)\\|_2^2
        &\le
        \Phi_A\big(k;\bar{\omega}_A,\beta_A\big)
        +
        \frac{8\Gamma_A(k)L\sigma_{\mathrm{sc}}^2d}{T\eta\tau N^2}
        \sum_{t=0}^{T-1}\frac1{b_t^2}
        \notag\\
        &\quad+
        \frac{16\Gamma_A(k)C_{\mathrm{dr}}}{T\eta\tau N^2}
        \sum_{t=0}^{T-1}
        \frac{\mathbb{E}[E_{\mathrm{clip},A}^t(k,b_t)]}{b_t^2}.
    \end{align}
\end{theorem}

\subsection{Construction of the Virtual Sequence and Separation}

Define the averaged error memory
\begin{equation}
    \bar e^{t-1}\triangleq \frac1N\sum_{i=1}^{N}e_i^{t-1}.
\end{equation}
The virtual sequence is defined by
\begin{equation}
    \tilde\theta^t\triangleq \theta^t+\bar e^{t-1},
    \qquad t\ge0,
\end{equation}
with $\bar e^{-1}=\mathbf{0}$. Equivalently, $\tilde\theta^{t+1}=\theta^{t+1}+\bar e^t$. Using~\eqref{eqn:transmit_signal} and the physical update~\eqref{eqn:physical_update_dr}, the error memories cancel out:
\begin{equation}
    \tilde\theta^{t+1}
    =\tilde\theta^t+a_A^t+e_{\mathrm{ch}}^t+e_{\mathrm{dr},A}^t,
\end{equation}
where
\begin{equation}
    a_A^t
    \triangleq
    \frac1N\sum_{i=1}^{N}\Delta_i^t-\bar\xi_A^t(k),
    \qquad
    \bar\xi_A^t(k)\triangleq\frac1N\sum_{i=1}^{N}\xi_{A,i}^t(k).
\end{equation}

By $L$-smoothness,
\begin{align}
    \label{eqn:descent_expansion}
    \mathbb{E}_t[f(\tilde\theta^{t+1})-f(\tilde\theta^t)]
    &\le
    \mathbb{E}_t\left\langle\nabla f(\tilde\theta^t),
    a_A^t+e_{\mathrm{ch}}^t+e_{\mathrm{dr},A}^t\right\rangle
    \notag\\
    &\quad+
    \frac{L}{2}\mathbb{E}_t\|a_A^t+e_{\mathrm{ch}}^t+e_{\mathrm{dr},A}^t\|_2^2.
\end{align}
Since $e_{\mathrm{ch}}^t$ is conditionally zero-mean, $\mathbb{E}_t\langle\nabla f(\tilde\theta^t),e_{\mathrm{ch}}^t\rangle=0$. No zero-mean or independence property is assumed for $e_{\mathrm{dr},A}^t$. We decompose
\begin{equation}
    \mathbb{E}_t[f(\tilde\theta^{t+1})-f(\tilde\theta^t)]
    \le T_{1,A}+T_{\mathrm{dr,lin},A}+T_{2,A},
\end{equation}
where
\begin{align}
    T_{1,A}&\triangleq\mathbb{E}_t\langle\nabla f(\tilde\theta^t),a_A^t\rangle,\\
    T_{\mathrm{dr,lin},A}&\triangleq\mathbb{E}_t\langle\nabla f(\tilde\theta^t),e_{\mathrm{dr},A}^t\rangle,\\
    T_{2,A}&\triangleq\frac{L}{2}\mathbb{E}_t\|a_A^t+e_{\mathrm{ch}}^t+e_{\mathrm{dr},A}^t\|_2^2.
\end{align}

\subsection{Bounding $T_{1,A}$, $T_{\mathrm{dr,lin},A}$, and $T_{2,A}$}

\subsubsection{Descent Directions $T_{1,A}$ and $T_{\mathrm{dr,lin},A}$}

Using $\Delta_i^t=-\eta\sum_{s=1}^{\tau}g_i^{t,s-1}$ and Assumption~\ref{assum:unbiased},
\begin{align}
    \mathbb{E}_t\left\langle\nabla f(\tilde\theta^t),\frac1N\sum_{i=1}^{N}\Delta_i^t\right\rangle
    &=-\eta\tau\mathbb{E}_t\left\langle\nabla f(\tilde\theta^t),
    \frac{1}{N\tau}\sum_{i,s}\nabla f_i(\theta_i^{t,s-1})\right\rangle.
\end{align}
Applying the identity $\langle a,b\rangle=\frac12\|a\|_2^2+\frac12\|b\|_2^2-\frac12\|a-b\|_2^2$ and $L$-smoothness gives
\begin{align}
    \mathbb{E}_t\left\langle\nabla f(\tilde\theta^t),\frac1N\sum_i\Delta_i^t\right\rangle
    &\le
    -\frac{\eta\tau}{2}\mathbb{E}_t\|\nabla f(\tilde\theta^t)\|_2^2
    \notag\\
    &\quad-
    \frac{\eta\tau}{2}\mathbb{E}_t\left\|\frac{1}{N\tau}\sum_{i,s}\nabla f_i(\theta_i^{t,s-1})\right\|_2^2
    \notag\\
    &\quad+
    \frac{\eta L^2}{2N}\sum_{i,s}\mathbb{E}_t\|\tilde\theta^t-\theta_i^{t,s-1}\|_2^2.
\end{align}
For the clipping-bias term, Young's inequality and Assumption~\ref{assum:clipping} yield
\begin{equation}
    -\mathbb{E}_t\langle\nabla f(\tilde\theta^t),\bar\xi_A^t(k)\rangle
    \le
    \frac{\eta\tau}{4}\mathbb{E}_t\|\nabla f(\tilde\theta^t)\|_2^2+
    \eta\tau\beta_A^2(k).
\end{equation}
Hence
\begin{align}
    \label{eqn:t1_bound_new}
    T_{1,A}
    &\le
    -\frac{\eta\tau}{4}\mathbb{E}_t\|\nabla f(\tilde\theta^t)\|_2^2
    +\frac{\eta L^2}{2N}\sum_{i,s}\mathbb{E}_t\|\tilde\theta^t-\theta_i^{t,s-1}\|_2^2
    \notag\\
    &\quad-
    \frac{\eta\tau}{2}\mathbb{E}_t\left\|\frac{1}{N\tau}\sum_{i,s}\nabla f_i(\theta_i^{t,s-1})\right\|_2^2
    +\eta\tau\beta_A^2(k).
\end{align}

For the linear dynamic-range term, Young's inequality and Lemma~\ref{lemma:dr_error} imply
\begin{align}
    \label{eqn:t_dr_lin_bound}
    T_{\mathrm{dr,lin},A}
    &\le
    \frac{\eta\tau}{16}\mathbb{E}_t\|\nabla f(\tilde\theta^t)\|_2^2
    +
    \frac{4}{\eta\tau b_t^2N^2}\mathbb{E}_t[E_{\mathrm{clip},A}^t(k,b_t)].
\end{align}

\subsubsection{Variance and Noise $T_{2,A}$}

Because $e_{\mathrm{dr},A}^t$ may be correlated with both the update and the channel noise, we separate it using Young's inequality. With a symmetric choice of the Young parameter,
\begin{align}
    \mathbb{E}_t\|a_A^t+e_{\mathrm{ch}}^t+e_{\mathrm{dr},A}^t\|_2^2
    &\le
    2\mathbb{E}_t\|a_A^t+e_{\mathrm{ch}}^t\|_2^2
    +2\mathbb{E}_t\|e_{\mathrm{dr},A}^t\|_2^2
    \notag\\
    &=2\mathbb{E}_t\|a_A^t\|_2^2
    +2\mathbb{E}_t\|e_{\mathrm{ch}}^t\|_2^2
    +2\mathbb{E}_t\|e_{\mathrm{dr},A}^t\|_2^2,
\end{align}
where the cross term between $a_A^t$ and $e_{\mathrm{ch}}^t$ vanishes by conditional zero-mean channel noise. Thus
\begin{equation}
    T_{2,A}
    \le
    L\mathbb{E}_t\|a_A^t\|_2^2
    +L\mathbb{E}_t\|e_{\mathrm{ch}}^t\|_2^2
    +L\mathbb{E}_t\|e_{\mathrm{dr},A}^t\|_2^2.
\end{equation}
Using Jensen's inequality and Assumption~\ref{assum:clipping},
\begin{equation}
    \mathbb{E}_t\|a_A^t\|_2^2
    \le
    \frac{2}{N}\sum_{i=1}^{N}\mathbb{E}_t\|\Delta_i^t\|_2^2
    +2\eta^2\tau^2\beta_A^2(k).
\end{equation}
The accumulated local update satisfies
\begin{align}
    \frac1N\sum_{i=1}^{N}\mathbb{E}_t\|\Delta_i^t\|_2^2
    &\le
    \eta^2\tau\sum_{s=1}^{\tau}
    \left(
    \frac{2L^2}{N}\sum_{i=1}^{N}\mathbb{E}_t\|\tilde\theta^t-\theta_i^{t,s-1}\|_2^2
    \right.
    \notag\\
    &\quad\left.
    +2G^2\mathbb{E}_t\|\nabla f(\tilde\theta^t)\|_2^2
    +2\kappa^2+\zeta^2
    \right).
\end{align}
Substituting this estimate, the channel-noise variance, and Lemma~\ref{lemma:dr_error}, we obtain
\begin{align}
    \label{eqn:t2_bound_new}
    T_{2,A}
    &\le
    4L\eta^2\tau^2G^2\mathbb{E}_t\|\nabla f(\tilde\theta^t)\|_2^2
    +\frac{4L^3\eta^2\tau}{N}\sum_{i,s}\mathbb{E}_t\|\tilde\theta^t-\theta_i^{t,s-1}\|_2^2
    \notag\\
    &\quad+
    2L\eta^2\tau^2(2\kappa^2+\zeta^2)
    +2L\eta^2\tau^2\beta_A^2(k)
    +\frac{Ld\sigma_{\mathrm{sc}}^2}{2b_t^2N^2}
    +\frac{L}{b_t^2N^2}\mathbb{E}_t[E_{\mathrm{clip},A}^t(k,b_t)].
\end{align}

\subsection{Bounding the Virtual Local Divergence}

Define
\begin{equation}
    D_{A,t}
    \triangleq
    \frac1N\sum_{i,s}\mathbb{E}_t\|\tilde\theta^t-\theta_i^{t,s-1}\|_2^2.
\end{equation}
Since $\tilde\theta^t-\theta_i^{t,s-1}=\bar e^{t-1}+\theta^t-\theta_i^{t,s-1}$, Lemma~\ref{lemma:local_div} and the gradient conversion
\begin{equation}
    \|\nabla f(\theta^t)\|_2^2
    \le
    2\|\nabla f(\tilde\theta^t)\|_2^2+2L^2\|\bar e^{t-1}\|_2^2
\end{equation}
give
\begin{align}
    \label{eqn:psi_bound}
    D_{A,t}
    &\le
    64\eta^2\tau^3G^2\mathbb{E}_t\|\nabla f(\tilde\theta^t)\|_2^2
    \notag\\
    &\quad+
    2\tau(1+32\eta^2\tau^2G^2L^2)\|\bar e^{t-1}\|_2^2
    +32\eta^2\tau^3\kappa^2+8\tau^2\eta^2\zeta^2.
\end{align}

\subsection{Learning Rate Conditions and Single-Step Descent}

Define the $\mathcal{F}_t$-measurable memory quantity
\begin{equation}
    \mathcal{M}_{A,t-1}\triangleq 2L^2\|\bar e^{t-1}\|_2^2.
\end{equation}
Combining~\eqref{eqn:t1_bound_new},~\eqref{eqn:t_dr_lin_bound}, and~\eqref{eqn:t2_bound_new}, substituting~\eqref{eqn:psi_bound}, and collecting terms gives
\begin{align}
    \mathbb{E}_t[f(\tilde\theta^{t+1})-f(\tilde\theta^t)]
    &\le
    -\frac{3\eta\tau}{16}\mathbb{E}_t\|\nabla f(\tilde\theta^t)\|_2^2
    +C_{\nabla,A}^{(0)}\mathbb{E}_t\|\nabla f(\tilde\theta^t)\|_2^2
    \notag\\
    &\quad+
    \Lambda\mathcal{M}_{A,t-1}
    +\eta^2\tau^2L\Psi
    +\mathcal{E}_A(k)
    \notag\\
    &\quad+
    \frac{Ld\sigma_{\mathrm{sc}}^2}{2b_t^2N^2}
    +C_{\mathrm{dr}}\frac{\mathbb{E}_t[E_{\mathrm{clip},A}^t(k,b_t)]}{b_t^2N^2},
\end{align}
where
\begin{equation}
    C_{\nabla,A}^{(0)}
    =4L\eta^2\tau^2G^2
    +64\eta^3\tau^3L^2G^2\left(\frac12+4L\eta\tau\right),
\end{equation}
\begin{equation}
    \Lambda
    =\eta\tau\left(\frac12+4L\eta\tau\right)
    \left(1+32\eta^2\tau^2G^2L^2\right),
\end{equation}
\begin{equation}
    \mathcal{E}_A(k)
    \triangleq
    (\eta\tau+2L\eta^2\tau^2)\beta_A^2(k),
    \qquad
    C_{\mathrm{dr}}\triangleq \frac{4}{\eta\tau}+L,
\end{equation}
and
\begin{equation}
    \Psi
    \triangleq
    2(2\kappa^2+\zeta^2)
    +\eta L\left(\frac12+4L\eta\tau\right)
    (32\tau\kappa^2+8\zeta^2),
\end{equation}
This constant is zero when $\kappa^2+\zeta^2=0$.

Under the learning-rate condition~\eqref{eqn:lr_condition_new}, one has $C_{\nabla,A}^{(0)}\le \eta\tau/16$. Therefore,
\begin{align}
    \label{eqn:single_step_descent_new}
    \mathbb{E}_t[f(\tilde\theta^{t+1})-f(\tilde\theta^t)]
    &\le
    -\frac{\eta\tau}{8}\mathbb{E}_t\|\nabla f(\tilde\theta^t)\|_2^2
    +\Lambda\mathcal{M}_{A,t-1}
    \notag\\
    &\quad+
    \eta^2\tau^2L\Psi
    +\mathcal{E}_A(k)
    +\frac{Ld\sigma_{\mathrm{sc}}^2}{2b_t^2N^2}
    \notag\\
    &\quad+
    C_{\mathrm{dr}}\frac{\mathbb{E}_t[E_{\mathrm{clip},A}^t(k,b_t)]}{b_t^2N^2}.
\end{align}

\subsection{Virtual Model Convergence}

Telescoping~\eqref{eqn:single_step_descent_new} over $t=0,\ldots,T-1$, using $f(\tilde\theta^T)\ge f^*$ and $\tilde\theta^0=\theta^0$, and defining $\Delta f\triangleq f(\theta^0)-f^*$, we obtain
\begin{equation}
    \label{eqn:virtual_bound_with_memory}
    \frac1T\sum_{t=0}^{T-1}\mathbb{E}\|\nabla f(\tilde\theta^t)\|_2^2
    \le
    B_A(k,\{b_t\})+
    \frac{8\Lambda}{T\eta\tau}\sum_{t=0}^{T-1}\mathbb{E}[\mathcal{M}_{A,t-1}],
\end{equation}
where
\begin{align}
    B_A(k,\{b_t\})
    &\triangleq
    \frac{8\Delta f}{T\eta\tau}
    +8\eta\tau L\Psi
    +\frac{8\mathcal{E}_A(k)}{\eta\tau}
    \notag\\
    &\quad+
    \frac{4L\sigma_{\mathrm{sc}}^2d}{T\eta\tau N^2}\sum_{t=0}^{T-1}\frac1{b_t^2}
    +
    \frac{8C_{\mathrm{dr}}}{T\eta\tau N^2}
    \sum_{t=0}^{T-1}
    \frac{\mathbb{E}[E_{\mathrm{clip},A}^t(k,b_t)]}{b_t^2}.
\end{align}


\subsection{Conversion to the Physical Model}

By $L$-smoothness,
\begin{equation}
    \mathbb{E}\|\nabla f(\theta^t)\|_2^2
    \le
    2\mathbb{E}\|\nabla f(\tilde\theta^t)\|_2^2+E_{A,t-1},
\end{equation}
where
\begin{equation}
    E_{A,t-1}
    \triangleq
    \mathbb{E}[\mathcal{M}_{A,t-1}]
    =2L^2\mathbb{E}\|\bar e^{t-1}\|_2^2.
\end{equation}
Using Jensen's inequality, $e_i^{-1}=\mathbf{0}$, and Lemma~\ref{lemma:error_bound},
\begin{align}
    \sum_{t=0}^{T-1}E_{A,t-1}
    &\le
    2L^2\mu_A(k)
    \sum_{t=0}^{T-1}\frac1N\sum_{i=1}^{N}\mathbb{E}\|\Delta_i^t\|_2^2.
    \label{eqn:memory_start}
\end{align}
It remains to close the recursion because the accumulated-update norm itself depends on the error memory through $\tilde\theta^t$.
From the accumulated-update estimate used in~\eqref{eqn:t2_bound_new},
\begin{align}
    \frac1N\sum_{i=1}^{N}\mathbb{E}\|\Delta_i^t\|_2^2
    &\le
    2\eta^2\tau L^2\mathbb{E}[D_{A,t}]
    +2\eta^2\tau^2G^2\mathbb{E}\|\nabla f(\tilde\theta^t)\|_2^2
    +\eta^2\tau^2(2\kappa^2+\zeta^2).
    \label{eqn:delta_accum_preclose}
\end{align}
Substituting the virtual local-divergence bound~\eqref{eqn:psi_bound} into~\eqref{eqn:delta_accum_preclose} gives
\begin{align}
    \frac1N\sum_{i=1}^{N}\mathbb{E}\|\Delta_i^t\|_2^2
    &\le
    c_{g}\,\mathbb{E}\|\nabla f(\tilde\theta^t)\|_2^2
    +c_{e}\,E_{A,t-1}
    +c_{0},
    \label{eqn:delta_accum_closed_raw}
\end{align}
where
\begin{align}
    c_g&=2\eta^2\tau^2G^2+128\eta^4\tau^4L^2G^2,\\
    c_e&=2\eta^2\tau^2\left(1+32\eta^2\tau^2G^2L^2\right),\\
    c_0&=\eta^2\tau^2(2\kappa^2+\zeta^2)
    +64\eta^4\tau^4L^2\kappa^2
    +16\eta^4\tau^3L^2\zeta^2.
\end{align}
Under~\eqref{eqn:lr_condition_new}, these coefficients satisfy the conservative bounds
\begin{equation}
    c_g\le 4\eta^2\tau^2G^2,
    \qquad
    c_0\le 2\eta^2\tau^2(2\kappa^2+\zeta^2),
    \qquad
    2L^2\mu_A(k)c_e\le \frac12.
    \label{eqn:memory_absorb_conditions}
\end{equation}
Combining~\eqref{eqn:memory_start}--\eqref{eqn:delta_accum_closed_raw} and using~\eqref{eqn:memory_absorb_conditions} yield
\begin{align}
    \sum_{t=0}^{T-1}E_{A,t-1}
    &\le
    2L^2\mu_A(k)c_g
    \sum_{t=0}^{T-1}\mathbb{E}\|\nabla f(\tilde\theta^t)\|_2^2
    +2L^2\mu_A(k)c_e
    \sum_{t=0}^{T-1}E_{A,t-1}
    +2L^2\mu_A(k)Tc_0
    \notag\\
    &\le
    8L^2\mu_A(k)\eta^2\tau^2G^2
    \sum_{t=0}^{T-1}\mathbb{E}\|\nabla f(\tilde\theta^t)\|_2^2
    +\frac12\sum_{t=0}^{T-1}E_{A,t-1}
    +4L^2\mu_A(k)T\eta^2\tau^2(2\kappa^2+\zeta^2).
\end{align}
Moving the middle term to the left gives the explicit memory-closure inequality
\begin{equation}
    \label{eqn:memory_sum_bound}
    \frac1T\sum_{t=0}^{T-1}E_{A,t-1}
    \le
    M_A(k)
    \frac1T\sum_{t=0}^{T-1}\mathbb{E}\|\nabla f(\tilde\theta^t)\|_2^2
    +R_A(k),
\end{equation}
where
\begin{equation}
    M_A(k)=16L^2\mu_A(k)\eta^2\tau^2G^2,
\end{equation}
and
\begin{equation}
    R_A(k)=8L^2\mu_A(k)\eta^2\tau^2(2\kappa^2+\zeta^2).
\end{equation}
The learning-rate condition also ensures
\begin{equation}
    \frac{8\Lambda M_A(k)}{\eta\tau}\le \frac12.
\end{equation}
Therefore, combining~\eqref{eqn:virtual_bound_with_memory} and~\eqref{eqn:memory_sum_bound} gives
\begin{equation}
    \frac1T\sum_{t=0}^{T-1}\mathbb{E}\|\nabla f(\tilde\theta^t)\|_2^2
    \le
    2B_A(k,\{b_t\})+\frac{16\Lambda R_A(k)}{\eta\tau}.
\end{equation}
Averaging the physical-to-virtual conversion over $T$ rounds and substituting~\eqref{eqn:memory_sum_bound}, we obtain
\begin{align}
    \frac1T\sum_{t=0}^{T-1}\mathbb{E}\|\nabla f(\theta^t)\|_2^2
    &\le
    \Gamma_A(k)
    \frac1T\sum_{t=0}^{T-1}\mathbb{E}\|\nabla f(\tilde\theta^t)\|_2^2
    +R_A(k),
\end{align}
where
\begin{equation}
    \Gamma_A(k)=2+M_A(k)=2+16L^2\mu_A(k)\eta^2\tau^2G^2.
\end{equation}
Substituting the bound on the virtual sequence yields
\begin{align}
    \label{eqn:final_dr_convergence_bound}
    \frac1T\sum_{t=0}^{T-1}\mathbb{E}\|\nabla f(\theta^t)\|_2^2
    &\le
    \frac{8\Gamma_A(k)L\sigma_{\mathrm{sc}}^2d}{T\eta\tau N^2}
    \sum_{t=0}^{T-1}\frac1{b_t^2}
    \notag\\
    &\quad+
    \frac{16\Gamma_A(k)C_{\mathrm{dr}}}{T\eta\tau N^2}
    \sum_{t=0}^{T-1}
    \frac{\mathbb{E}[E_{\mathrm{clip},A}^t(k,b_t)]}{b_t^2}
    +\Phi_A\big(k;\bar{\omega}_A,\beta_A\big),
\end{align}
where $\Phi_A\big(k;\bar{\omega}_A,\beta_A\big)$ is defined in~\eqref{eqn:phi_def_theorem}. This proves Theorem~\ref{thm:main_convergence}.

\begin{corollary}[\bf Joint Utility and Per-Round Privacy]
    \label{cor:joint_utility_privacy}
    Under Theorem~\ref{thm:main_convergence}, use $b_t=b_t^\star(k)$ from~\eqref{eqn:per_round_scaling}. Then the physical iterates satisfy the convergence bound, and every round individually satisfies replace-one client-level $(\epsilon,\delta)$-DP. No cross-round privacy-budget allocation is required.
\end{corollary}
\begin{proof}
    Convergence follows from Theorem~\ref{thm:main_convergence} and Assumption~\ref{assum:public_scaling}. Per-round privacy follows from Theorem~\ref{thm:per_round_privacy}.
\end{proof}

\subsection{Interpretation and Dynamic-Range-Surrogate Interface}

The final convergence bound has three groups. The learning-side term collects local drift, stochastic-gradient variance, heterogeneity, error-feedback memory, and pre-transmission clipping bias. The recovered channel-noise term is proportional to
\begin{equation}
    \frac1T\sum_{t=0}^{T-1}\frac{1}{b_t^2N^2},
\end{equation}
and the receiver dynamic-range term is proportional to
\begin{equation}
    \frac1T\sum_{t=0}^{T-1}
    \frac{\mathbb E[E_{\mathrm{clip},A}^t(k,b_t)]}{b_t^2N^2}.
\end{equation}
There is no client-participation sampling term because every logical round aggregates all $N$ clients. Path loss and fading affect convergence only through the public per-round ceiling $B_P^t(k)$, the resulting $b_t^\star(k)$, and the realized receiver waveform.

The per-round privacy ceiling and power ceiling both scale as $1/\sqrt{k}$ under the fixed coordinate-clipping rule. Hence sparse transmission can increase the feasible $b_t^\star(k)$ relative to full transmission and reduce recovered noise. At equal $k$, Top-$k$ and Rand-$k$ have the same worst-case replace-one sensitivity ceiling; their distinction is represented by the certified retention coefficient, clipping bias, and support-overlap-dependent waveform residual.

For calibrated operating-point selection, the exact radial term may be summarized by
\begin{equation}
    \widehat D_{\mathrm{DR},A}(k;\gamma)
    =\frac{1}{T_{\mathrm{cal}}SQ}
    \sum_{t=1}^{T_{\mathrm{cal}}}\sum_{\ell=1}^{S}\sum_{n=0}^{Q-1}
    \left(|\bar y_{A,\ell}^t[n;k]|-\gamma\right)_+^2,
\end{equation}
where $\bar y_{A,\ell}^t=a_ty_{A,\ell}^t$ is the RMS-normalized noisy waveform. This is a design surrogate, not a replacement for the exact theorem term. With frozen normalized dynamic-range statistics, the per-round surrogate is decreasing in $b_t$ through the recovered-noise term, so the largest feasible value is precisely
\begin{equation}
    b_t^\star(k)=\min\{B_{\epsilon,\mathrm{ex}}(k),B_P^t(k)\}.
\end{equation}
The exact nonlinear waveform problem is not claimed to be globally optimized by this closed form.

\end{document}
